At least one predeclared multiplicity-controlled contrast separated from zero. The result supports only the observed internal operational-target comparison and does not validate the optical residual as process physics. The generator diagnostics reinforce the distinction between numerical regularization and downstream utility: the validation reconstruction path, morphology distances, and detector endpoint address different questions and need not rank methods identically. The donor-matched non-neural and additional-real-exposure arms are particularly important because gains over a weak real-only baseline can otherwise be attributed to target reuse or additional optimizer exposure.

Three limitations set the boundary of interpretation. First, the inherited box is an operational region rather than an independently adjudicated particle label; no claim is made about particle count, streak identity, or thermophysical state. Second, the 176-specimen corpus provides a specimen-grouped internal hold-out but no independent machine, build, site, material, or acquisition campaign, so external validity remains unknown. Third, the coordinate field operates at 64 pixels and the higher-resolution branch was unsupported by native information; the resolution audit quantifies a pathway effect without recovering absent detail.

Within those limits, the study contributes a falsifiable controlled evaluation rather than a physics-validity claim. Future work should obtain independently adjudicated physical targets and a prospectively acquired external cohort before using the residual or detector for process inference. A calibrated time-resolved measurement would also be required before replacing the optical operator with a governing thermal or flow residual.
